좌표 압축(coordinate compression) 또는 값 압축(value compression)이라고 불리는 테크닉은 전체 값들의 절대적인 값보다 상대적인 대소관계만이 영향을 주는 상황에서 값의 범위를 줄이는 방법이다. 입력의 개수에 비해 값의 범위가 매우 넓을 때 유용하게 사용할 수 있다. 특히 값 자체를 인덱스의 의미로 사용해야 하는 경우에 유용하다. 이번 장에서는 좌표 압축을 하는 방법과 사용 예시를 배운다.

\section{값의 범위를 줄이기}
다음 \Cref{ex: value-compression}을 보자.
\begin{example} \label{ex: value-compression}
    이차원 좌표평면에 있는 $N$개의 점이 주어진다. 각 좌표 $(x_i, y_i)$에 대해 $x$좌표가 더 작은 점의 수를 출력하는 프로그램을 작성하라. 즉, 총 $N$개의 수를 출력해야 한다.
\end{example}

가장 단순한 방법은 $N$개를 전부 보면서 $(x_i, y_i)$가 등장하면 $x_i$의 개수에 $1$을 더한 이후 누적 합을 사용하는 것이다. 하지만 $x_i$의 범위가 매우 넓으면 문제가 된다. 만약 $x_i$가 $10^9$ 이하의 자연수라면 $10^9$ 크기의 배열을 만들어야 할 것이다.

그런데 굳이 $10^9$ 이하의 모든 값에 대한 정보를 저장하는 배열을 만들어야 할까? 최대 $N$개의 값이 들어오기 때문에, 개수를 세는 배열에서 최대 $N$개의 값만 $0$이 아니고, $0$이 아닌 인덱스만 고려해도 충분하다. 이러한 관찰에서 시작된 알고리즘이 좌표 압축이다. 좌표 압축의 정의는 어떤 배열 $A$에서 각 $A_i$를 $A_i$ 이하의 수의 종류의 수로 정하는 것이다. 간단하게 말해, $A_i$가 배열 $A$에서 얼마나 큰지를 상대적으로 나타내는 값이다. $A_i$가 $A$에서 $k$번째로 작으면 좌표 압축한 배열에서 $A_i$는 $k$이다.

이렇게 좌표 압축을 하면 빠지는 점 없이 모든 점을 처리할 수 있으며, 각 $j$에 대해 $A_i=j$인 $i$가 적어도 하나는 존재하기 때문에 기존의 목표를 달성하였다고 볼 수 있다. 이제 마지막으로 남은 문제는 좌표 압축을 한 배열에서 답을 구해도 되는지이다. 우리는 다음 관찰 하나로 좌표 압축된 배열에서 답을 구해도 된다는 사실을 정당화할 수 있다.

\begin{obs}
    좌표 압축한 이후에 수의 대소관계는 바뀌지 않는다. 다른 말로 원래 배열에서 더 컸던 수는 좌표 압축한 배열에서도 여전히 더 크다.
\end{obs}

문제에서 $x_i$보다 $x$좌표가 작은 점의 개수를 찾으라고 하였으므로 좌표 압축을 한 배열에서 문제를 풀어도 전혀 문제가 없다.

마지막으로 위 내용을 수학적으로 명확하게 정리해보자. $A$를 좌표 압축한 배열을 $B$라고 할 때, 
\begin{itemize}
    \item $B$의 모든 원소는 $A$의 길이 이하이다.
    \item $A_i < A_j$와 $B_i < B_j$는 필요충분관계, $A_i = A_j$와 $B_i = B_j$도 필요충분관계이다.
\end{itemize}

그렇다면 좌표 압축은 어떻게 할까? 다음 \Cref{ex: value-compression}에 대한 정답 코드인 \Cref{code: value-compression}로 알아보자.
\begin{code} \label{code: value-compression}
\Cref{ex: value-compression}의 정답 코드
\begin{lstlisting}
#include <bits/stdc++.h>
using namespace std;

int n, cnt[101010], sum[101010];
pair<int, int> arr[101010];
vector<int> val;

int main() {
    cin >> n;
    for (int i = 1; i <= n; i++) cin >> arr[i].first >> arr[i].second;

    for (int i = 1; i <= n; i++) val.push_back(arr[i].first);
    sort(val.begin(), val.end());
    val.erase(unique(val.begin(), val.end()), val.end());
    for (int i = 1; i <= n; i++) 
        arr[i].first = lower_bound(val.begin(), val.end(), arr[i].first) - val.begin() + 1;

    for (int i = 1; i <= n; i++) cnt[arr[i].first]++;
    for (int i = 1; i <= n; i++) sum[i] = sum[i - 1] + cnt[i];
    for (int i = 1; i <= n; i++) cout << sum[arr[i].first - 1] << " ";

    return 0;
}
\end{lstlisting}
\end{code}
위 코드에서 좌표 압축에 해당하는 부분은 13, 14, 16번 줄이다. 세 줄에 적힌 코드는 각각 다음과 같은 역할을 한다.
\begin{itemize}
    \item 13, 14번째 줄은 \texttt{val}를 정렬하고 중복된 값을 제거하는 코드이다.
    \item 16번째 줄은 \texttt{arr}에 포함된 서로 다른 값을 바탕으로 좌표 압축한 새로운 배열을 만드는 코드이다.
\end{itemize}
13번째 줄에 있는 \texttt{std::unqiue}에 대한 자세한 내용은 다음 C++ 레퍼런스 사이트에서 확인할 수 있다: \url{https://en.cppreference.com/w/cpp/algorithm/unique}

모든 좌표 압축은 코드가 동일하다. 그리고 모든 좌표 압축은 \textbf{코드 세 줄(13, 14, 16번 줄)}로 끝난다. 매우 자주 쓰이는 테크닉이니 꼭 이 코드 세 줄 만큼은 외워두도록 하자.

\section{연습문제 A}
\begin{problem}
    \Cref{ex: value-compression}을 좌표 압축 없이 해결하는 코드를 작성하라.
\end{problem}
\begin{problem}
    \Cref{code: value-compression}에서 다시 원래 배열의 값을 알고 싶다면 어떻게 하면 되겠는가?
\end{problem}

\begin{problem}
    (18870번) 주어진 길이 $N$의 수열에 좌표압축을 한 결과를 출력하라. 시작은 $0$번이며, $N \leq 10^6$, 수열의 값은 절댓값이 $10^9$ 이하인 정수이다.
\end{problem}

\begin{problem}
    (18869번) 길이 $N$인 $M$개의 수열이 주어질 때, $\binom{M}{2}$개의 쌍들 중 다음 조건을 만족하는 조합의 개수를 출력하라.
    \begin{itemize}
        \item 수열 $A_1, A_2, \cdots, A_N$과 $B_1, B_2, \cdots, B_N$이 임의의 $1 \leq i, j \leq N$에 대하여 
        \begin{align*}
            A_i < A_j &\Rightarrow B_i < B_j \\
            A_i = A_j &\Rightarrow B_i = B_j \\
            A_i > A_j &\Rightarrow B_i > B_j
        \end{align*}
        를 만족해야 한다.
    \end{itemize}
    $2 \leq M \leq 100$, $3 \leq N \leq 10,000$이며 수열의 값은 $10^6$ 이하의 자연수이다.
\end{problem}

\section{연습문제 B}

\begin{problem}
    (2370번) $n$개의 구간 $[s_i, e_i]$가 주어지면 작은 $i$부터 $i$번 구간 $[s_i, e_i]$을 $i$번 색으로 칠하자. 이미 색이 칠해져 있다면 새로운 색으로 덮어서 칠하는 것이다. $n$번 색까지 칠한 이후에 보이는 색의 종류의 수를 출력하라. $1 \leq n \leq 10^4$, $1 \leq s_i \leq e_i \leq 10^8$이다.
\end{problem}

\begin{problem} 
    (5926번) 수직선 상의 $n$개의 점이 있고, 각 점에는 숫자가 적혀 있다. 다음 조건을 만족하도록 만들 수 있는 구간 $[s, e]$를 생각하자.
    \begin{itemize}
        \item 임의의 숫자 $k$에 대해 $n$개의 점들 중 $k$가 적힌 점이 존재한다면, $[s, e]$에 $k$에 적힌 점이 존재해야 한다.
    \end{itemize}
    가능한 $e - s$의 최솟값을 출력하라. $1 \leq n \leq 5 \times 10^4$이고 점의 좌표와 점에 적힌 숫자는 모두 $10^9$ 이하의 양의 정수이다.
\end{problem}

\begin{problem}
    (11540번) $N$개의 문제 중 플레이어 A가 풀 수 있는 $A$개의 문제 번호와 플레이어 B가 풀 수 있는 $B$개의 문제 번호가 주어진다. 그들은 $1$번 문제부터 포면서 두 플레이어 중 한 명 이상이 풀 수 있다면 풀 것이다. 이때 문제를 풀 사람이 컴퓨터를 사용해야 한다. 이렇게 하여 최대한 많은 문제를 풀 때 컴퓨터를 사용하는 사람을 교체해야 하는 횟수를 출력하라. 처음 컴퓨터를 사용할 사람은 자유롭게 정할 수 있다. $1 \leq N \leq 10^9$, $1 \leq A, B \leq 5 \times 10^4$이다.
\end{problem}

\begin{problem}
    (11997번) 이차원 좌표평면에 $n$개의 점이 있다. 이 좌표평면을 짝수 $a, b$에 대해 $x=a$와 $y=b$로 네 부분으로 나누었을 때, 네 부분에 있는 점의 개수의 최댓값의 가능한 최솟값을 출력하라. $1 \leq n \leq 1000$이고 점의 $x$좌표와 $y$좌표는 모두 $10^6$ 이하의 양의 홀수이다.
\end{problem}

\begin{problem}
    (16806번) 좌표평면에 $n$개의 점이 있고, 각 점은 상하좌우 중 하나의 \textit{고유 방향}을 가진다. 이제 좌표평면 상의 원하는 위치에 상하좌우 중 한 방향으로 공을 굴리자. 그러면, 공이 점을 지나는 순간 공의 방향은 지난 점의 \textit{고유 방향}으로 바뀌고, 점은 사라진다. 이렇게 계속 운동하면 공이 유한 개의 점만 지남을 보장할 수 있고, 지난 점의 개수가 점수가 된다. 점수가 최대가 되도록 공의 위치와 이동 방향을 잘 정할 때, 가능한 최대 점수를 출력하라. $1 \leq n \leq 3000$이고 점의 $x$좌표와 $y$좌표는 모두 절댓값이 $10^9$ 이하인 정수이다.
\end{problem}

\begin{problem}
    (18205번) \textbf{소숫점 $6$ 자리}로 주어지는 $n$개의 수와 $m$개의 쿼리가 있다. 쿼리의 내용은 다음과 같다:
    \begin{itemize}
        \item $1 \leq l \leq r \leq n$이 주어질 때, $x = A_i$인 $l \leq i \leq r$인 $i$가 $\left\lfloor \frac{r-l+1}{2} \right\rfloor + 1$개 이상(과반)인 $x$가 존재하는지 판별하라.
    \end{itemize}
    각 쿼리에 대해 쿼리가 참이면 \texttt{usable}을, 거짓이면 \texttt{unusable}을 출력하라. $1 \leq n \leq 10^4$, $1 \leq m \leq 10^4$이고 주어지는 수는 $0.000000$ 이상 $14.000000$ 이하이다.
\end{problem}

\begin{problem}
    (18707번) 각 칸마다 숫자가 적혀있는 $N \times M$ 크기의 격자는 $K$개의 칸에는 $1$이 적혀있고 나머지 칸에는 전부 $0$이 적혀 있다. $[1, i] \times [1, j]$의 직사각형 영역(prefix sub-grid)에 적혀 있는 $1$의 개수가 홀수개인 $(i, j)~(1 \leq i \leq N, 1 \leq j \leq M)$의 개수를 구하여라. \textbf{여러 개의 테스트케이스에 대해 해결해야 함에 유의하라.} $1 \leq N, M \leq 10^9$, $0 \leq K \leq 10^3$이고, 입력으로는 $1$이 적혀 있는 칸의 좌표가 주어진다.
\end{problem}

\begin{problem} 
    (20440번) $n$개의 구간 $[s_i, e_i)$가 주어질 때, $x \in [s_i, e_i)$인 $i$의 개수의 최댓값과, 최대가 될 때 $x$의 값의 범위를 구간 $[s, e)$으로 표현하라. 단, 둘 이상이 가능하다면, \textbf{가장 시작 좌표가 작은 구간 하나}만 출력하라. $1 \leq n \leq 10^6$, $0 \leq s_i \leq e_i \leq 2.1 \times 10^9$이다.
\end{problem}

\begin{problem}
    (24572번) $N \times N$ 크기의 격자에 $T$개의 나무가 있다. 격자에 포함되면서 나무를 포함하지 않는 정사각형 격자의 한 변의 길이로 가능한 최댓값을 출력하라. $N \leq 5 \times 10^5$, $T \leq 100$이다.
\end{problem}

\section{연습문제 C}

\begin{problem}
    (2361번) 이차원 좌표평면에 $n$개의 점이 있다. 이 좌표평면에서 $\left\lceil \frac{n}{2} \right\rceil$개 이상(과반)의 점을 포함하는 최소 크기의 정사각형의 한 변의 길이를 찾아라. 정사각형의 경계에 있는 것도 정사각형에 포함되는 것으로 간주한다. $4 \leq n \leq 1000$이고 점의 $x$좌표와 $y$좌표는 모두 $100,015$ 이하의 \textbf{양의 실수로 소숫점 $6$자리}까지 주어질 수 있다.
\end{problem}

\begin{problem}
    (3324번) $W \times H$ 크기의 격자에 $N$개의 칸이 흰색 또는 검은색으로 칠해져 있고, 칠해진 칸에는 숫자가 적혀 있다.
    \begin{itemize}
        \item 두 격자 사이의 거리를 상하좌우로 인접한 칸으로 이동할 때 필요한 이동 횟수로 정의하자. 즉, $(a, b)$ 위치와 $(c, d)$ 위치 사이의 거리는 $|a-c| + |b-d|$이다. 
        \item 흰색 또는 검은색으로 칠해진 칸과 그 칸에 적혀 있는 숫자 $d$에 대하여 거리가 $d$ 이하인 점은 도달할 수 있다.
    \end{itemize}
    각 칸에 도달할 수 있는 흰색 칸의 수가 검은색 칸의 수보다 많은 칸의 개수와, 검은색 칸의 수가 흰색 칸의 수보다 많은 칸의 개수를 순서대로 출력하라. $1 \leq W, H \leq 10^9$, $1 \leq N \leq 3000$이고 각 칠해진 칸에 적혀 있는 수 $d$는 $0 \leq d < 5 \times 10^8$이다.
\end{problem}

\begin{problem}
    다음과 같은 방식으로 진행되는 투표가 있다:
    \begin{itemize}
        \item 후보는 두 명이어야만 한다. 선착순으로(?!) 두 명만 받는다.
        \item 각 유권자는 최대 $100,000$개의 표를 두 후보에게 나누어 던질 수 있다. 이때 한 표도 행사하지 않는 것은 불가능하다. 즉, 1번 후보에게 $a~(\geq 0)$개의 표를, 2번 후보에게 $b~(\geq 0)$개의 표를 던졌다면 $1 \leq a+b \leq 100,000$이어야 한다.
        \item 두 후보가 받은 표의 합이 더 큰 후보가 당선된다. 만약 같다면 두 후보를 공동 당선(?!) 처리한다.
    \end{itemize}
    여론조사에서는 2번 후보가 압도적 지지를 받았지만, 암흑의 단체에서 해킹을 통해 다음과 같은 방식의 로직을 추가하여 1번 후보만 당선되도록 조작하였다.
    \begin{itemize}
        \item 투표 결과의 전산 집계에서, 한 후보에게 $x$표 이상을 투표했거나, 전체 $y$표 이상을 투표한 경우 무시한다. 다른 말로, $a \geq x$, 또는 $b \geq x$, 또는 $a+b \geq y$인 표는 무시한다.
    \end{itemize}
    여러분은 공직선거법 위반 혐의로 암흑의 단체를 기소하기 위해 $1 \leq x, y \leq 10^5$를 만족하는 모든 $(x, y)$ 쌍을 찾아 검토해야 한다. $n$명의 유권자가 1번 후보와 2번 후보에게 던진 표의 수가 주어질 때, 가능한 $(x, y)$ 쌍의 경우의 수를 구하는 프로그램을 작성하라. $1 \leq n \leq 1000$이다.
\end{problem}

\begin{problem} \url{https://www.acmicpc.net/problem/28461} \end{problem}